\chapter{Экспериментальные исследования системы управления СП} \label{chapt6}

\section{Исследование точности позиционирования платформы} \label{sect6_1}

Для проверки конструктивной безопасности позиционирования в установившемся режиме были проведены испытания по измерению рабочей зоны 3DOF СП, характеристический портрет которой приведен на рисунке \ref{img:experimentStaticAccuracy}

\begin{figure}
\center
\includegraphics[width=0.95\linewidth]{InklinometerH0}
\caption{Эксперимент по проверке точности позиционирования СП}
\end{figure}

\begin{figure}[!htbp]
    \center{
        \subcaptionbox{\label{img:InklinometerHm}} {\includegraphics[width=0.95\linewidth]{InklinometerHm}}
        \hfill
        \subcaptionbox{\label{img:InklinometerHp}} {\includegraphics[width=0.95\linewidth]{InklinometerHp}}
	}
	
	1 - заданное значение, 2 - инклинометр
    \caption{Эксперимент по проверке точности позиционирования СП} % Этот текст попадает в названия рисунков в списке рисунков
    \label{img:experimentStaticAccuracy}
\end{figure}

На рисунке \ref{img:experimentStaticAccuracy} изображены срезы рабочей зоны СП в плоскости крен-тангаж. В ходе измерений производилось позиционирование платформы по траектории, показанной сплошной линией (Заданное положение), с одновременным измерением углов наклона инклинометром с погрешностью $0,5^o$. Как видно из графика, задаваемая траектория сходится с действительной в пределах погрешности прибора.

\section{Эксперимент по параметрической идентификации модели СП} \label{sect6_2}
Эксперимент по идентификации параметров модели производился следующим образом: 
\begin{itemize}

\item	платформа позиционировалась раздельно по координатным осям в пределах характеристики <<крен-тангаж>>, приведенной на рисунке \ref{img:experimentStaticAccuracy} с шагом $2^o$. На каждом шаге в течение выбранного интервала времени (~$60 c$) записывались значения вращающих моментов двигателей, рассчитываемых встроенными контроллерами преобразователей частоты. Определялись средние значения моментов на каждом шаге.
\item	Во время эксперимента измерялись электромагнитные моменты двигателей, т.к. непосредственное определение вращающих моментов на валах мотор-редукторов затруднительно. Расчет электромагнитных моментов производился средствами встроенного ПО преобразователей частоты на основе непосредственного датчиков токов фаз и текущих скоростей двигателей по известным для АД соотношениям для векторного управления моментом \cite{2014a}. На основе полученных точек $(q_i, G_i)_j: i = 1 \ldots 3, j = 1 \ldots n$, где $n$ – количество измерений, производилась идентификация параметров модели в соответствии с описанным выше алгоритмом
\item	далее производилось движение по синусоиде с одновременным считыванием вращающих моментов. По полученным данным проверялось расхождение измеренных и рассчитанных по модели значений вращающих моментов мотор-редукторов.
\end{itemize}
Результаты первой части эксперимента приведены в табл. 1.
Данные модели в таблице \ref{tbl:ModelVerificationExperiment} получены на основе идентификации, в результате которой определены следующие значения параметров:

\begin{table} [!htbp]%
    \centering
	\caption{Эксперимент по верификации динамической модели}%
	\label{tbl:ModelVerificationExperiment}
    \renewcommand{\arraystretch}{1.5}
    \begin{SingleSpace}
	\begin{tabulary}{\textwidth}{@{}>{\zz}L >{\zz}L >{\zz}L >{\zz}L >{\zz}L >{\zz}L >{\zz}L >{\zz}L >{\zz}L@{}}
        \toprule
		$q_1, ^o$	& $q_2, ^o$	& $q_3, ^o$	& $\bar{Q_1}$	& $\bar{Q_2}$	& $\bar{Q_3}$	& $Q_1$	& $Q_2$	& $Q_3$ \\
		\midrule
		150,0	& 150,0	& 60,0	& 68,0	& 114,1	& 489,5	& 78,6	& 114,6	& 521,5 \\
		150,0	& 150,0	& 90,0	& 115,9	& 154,9	& 506,7	& 153,0	& 152,2	& 602,8 \\
		150,0	& 150,0	& 120,0	& 169,9	& 195,6	& 394,7	& 201,7	& 231,0	& 487,5 \\
		150,0	& 150,0	& 150,0	& 203,5	& 236,4	& 191,4	& 228,7	& 313,8	& 244,7 \\
		97,7	& 97,7	& 97,7	& 269,0	& 297,0	& 244,5	& 284,0	& 305,0	& 250,4 \\
		108,2	& 87,1	& 97,7	& 205,5	& 360,1	& 266,8	& 218,0	& 355,8	& 250,8 \\
		118,9	& 75,5	& 97,7	& 112,5	& 398,7	& 252,7	& 139,5	& 390,5	& 250,7 \\
		130,3	& 62,6	& 97,7	& 26,2	& 412,0	& 252,7	& 52,6	& 403,6	& 249,5 \\
		143,3	& 47,1	& 97,7	& -32,5	& 351,6	& 236,3	& -34,1	& 383,5	& 246,0 \\
		93,3	& 93,3	& 112,2	& 318,6	& 342,4	& 176,8	& 324,6	& 338,3	& 185,8 \\
		88,8	& 88,8	& 127,3	& 346,2	& 389,3	& 94,7	& 364,2	& 371,7	& 105,4 \\
		84,3	& 84,3	& 145,2	& 362,7	& 423,6	& 27,2	& 407,9	& 401,2	& 25,0 \\
		\bottomrule
	\end{tabulary}%
   	\end{SingleSpace}
\end{table}

\begin{figure}[!htbp]
    \center{
        \subcaptionbox[List-of-Figures entry]{\label{img:IdentConstantAngles}} {\includegraphics[width=0.9\linewidth]{IdentConstantAngles}}%
		\hfill        
        \subcaptionbox{\label{img:IdentConstantTorques}} {\includegraphics[width=0.9\linewidth]{IdentConstantTorques}}
    }

    Тестирование модели СП в статикe
    \caption{Эксперимент по идентификации СП} % Этот текст попадает в названия рисунков в списке рисунков
    \label{img:identificationExperimentChartsConstant}
\end{figure}

\begin{figure}[!htbp]
    \center{
        \subcaptionbox[List-of-Figures entry]{\label{img:IdentSinesAngles}} {\includegraphics[width=0.9\linewidth]{IdentSinesAngles}}%
		\hfill        
        \subcaptionbox{\label{img:IdentSinesTorques}} {\includegraphics[width=0.9\linewidth]{IdentSinesTorques}}
    }

    Тестирование модели СП в статикe
    \caption{Эксперимент по идентификации СП} % Этот текст попадает в названия рисунков в списке рисунков
    \label{img:identificationExperimentChartsSines}
\end{figure}

Пунктирными линиями показаны измеренные значения моментов, сплошными – расчетные. На рисунке \ref{img:identificationExperimentChartsSines} показаны результаты исследования динамических процессов. Максимальное запаздывание расчетных данных по фазе относительно экспериментальных составило порядка 40 градусов. 
Расхождения между моделью и экспериментом можно объяснить задержками в канале передачи данных между частотными преобразователями и контроллером движения, а также особенностями работы внутренних контуров регулирования АД.
В дальнейших исследованиях предполагается использовать более высокоскоростной интерфейс передачи данных и изменить схему проведения эксперимента.

\section{Эксперимент по отработке типовых воздействий на системе подвижности} \label{sect6_3}

На рисунке \ref{img:testEKGturn} показаны результаты испытаний платформ с приводами вращательного типа. 
Для реализации системы управления использовался подход, примененный в системе 1. 
По графикам отработки задающих воздействий в осях крен-тангаж видно, что процессы при реакции на ступенчатые сигналы (рисунок \ref{img:Velocities_LFpulse}) апериодические, 
время переходного процесса определяется предельной скоростью платформы и транспортной задержкой. 
При отработке синусоидального сигнала (рисунок \ref{img:testLowFrequencySines}) с увеличением частоты возрастают перемещения по неуправляемым осям. 
Для улучшения качества переходных процессов применены алгоритмы КД системы 2, позволившие устранить паразитные движения при отработке типовых воздействий.

\begin{figure}[h!]
    \center{
        \subcaptionbox[List-of-Figures entry]{\label{img:RollPitch_EKGturn}} {\includegraphics[width=0.9\linewidth]{RollPitch_EKGturn}}%
        \hfill
        \subcaptionbox{\label{img:Velocities_EKGturn}} {\includegraphics[width=0.9\linewidth]{Velocities_EKGturn}}
    }

    На рисунке а) показаны текущие (2,4) и заданные (1,3) крен и тангаж платформы в градусах, на рисунке б) -- скорости первого (1), второго (2) и третьего (3) сервоприводов в об/мин. 
    \caption{Отработка эффекта наклона ЭКГ-18Р при поворотах} % Этот текст попадает в названия рисунков в списке рисунков
    \label{img:testEKGturn}
\end{figure}

\begin{figure}[h!]
    \center{
        \subcaptionbox[List-of-Figures entry]{\label{img:Height_EKG_digging}} {\includegraphics[width=0.9\linewidth]{Height_EKG_digging}}%
        \hfill
        \subcaptionbox{\label{img:Velocities_EKG_digging}} {\includegraphics[width=0.9\linewidth]{Velocities_EKG_digging}}
    }

    На рисунке а) показано текущее (2) и заданное (1) вертикальное смещение платформы в см, на рисунке б) -- скорости первого (1), второго (2) и третьего (3) сервоприводов в об/мин. 
    \caption{Отработка вибрации ЭКГ-18Р при копании} % Этот текст попадает в названия рисунков в списке рисунков
    \label{img:testEKGdigging}
\end{figure}

\begin{figure}[h!]
    \center{
        \subcaptionbox[List-of-Figures entry]{\label{img:Height_EKG_move}} {\includegraphics[width=0.9\linewidth]{Height_EKG_move}}%
        \hfill
        \subcaptionbox{\label{img:Velocities_EKG_move}} {\includegraphics[width=0.9\linewidth]{Velocities_EKG_move}}
    }

    На рисунке а) показано текущее (2) и заданное (1) вертикальное смещение платформы в см, на рисунке б) -- скорости первого (1), второго (2) и третьего (3) сервоприводов в об/мин. 
    \caption{Отработка вибрации ЭКГ-18Р при движении по забою} % Этот текст попадает в названия рисунков в списке рисунков
    \label{img:testEKGmovement}
\end{figure}

\begin{figure}[h!]
    \center{
        \subcaptionbox[List-of-Figures entry]{\label{img:RollPitch_lowfreq}} {\includegraphics[width=0.9\linewidth]{RollPitch_lowfreq}}%
        \hfill
        \subcaptionbox{\label{img:Velocities_lowfreq}} {\includegraphics[width=0.9\linewidth]{Velocities_lowfreq}}
    }
    На рисунке а) показаны текущие (2,4) и заданные (1,3) крен и тангаж платформы в градусах, текущее(6) и заданное (5) вертикальное смещение в см, на рисунке б) -- скорости первого (1), второго (2) и третьего (3) сервоприводов в об/мин. 
    \caption{Отработка низкочастотных синусоидальных колебаний} % Этот текст попадает в названия рисунков в списке рисунков
    \label{img:testLowFrequencySines}
\end{figure}
    
\begin{figure}[h!]
    \center{
        \subcaptionbox[List-of-Figures entry]{\label{img:RollPitch_LFlines}} {\includegraphics[width=0.9\linewidth]{RollPitch_LFlines}}%
        \hfill
        \subcaptionbox{\label{img:Velocities_LFlines}} {\includegraphics[width=0.9\linewidth]{Velocities_LFlines}}
    }
    На рисунке а) показаны текущие (2,4) и заданные (1,3) крен и тангаж платформы в градусах, текущее(6) и заданное (5) вертикальное смещение в см, на рисунке б) -- скорости первого (1), второго (2) и третьего (3) сервоприводов в об/мин.
    \caption{Отработка низкочастотных пилообразных колебаний} % Этот текст попадает в названия рисунков в списке рисунков
    \label{img:testLFlines}
\end{figure}
       
\begin{figure}[h!]
    \center{
        \subcaptionbox[List-of-Figures entry]{\label{img:RollPitch_LFpulse}} {\includegraphics[width=0.9\linewidth]{RollPitch_LFpulse}}%
        \hfill
        \subcaptionbox{\label{img:Velocities_LFpulse}} {\includegraphics[width=0.9\linewidth]{Velocities_LFpulse}}
    }
    На рисунке а) показаны текущие (2,4) и заданные (1,3) крен и тангаж платформы в градусах, текущее(6) и заданное (5) вертикальное смещение в мм, на рисунке б) -- скорости первого (1), второго (2) и третьего (3) сервоприводов в об/мин.
    \caption{Отработка импульсов по крену и тангажу с одновременной вибрацией по высоте} % Этот текст попадает в названия рисунков в списке рисунков
    \label{img:testPulsesRollPitch}
\end{figure}
\clearpage
